
\newpage

\vspace*{\fill}

THIS PAGE INTENTIONALLY LEFT BLANK

I WOULDN'T LEAVE THIS BLANK UNINTENTIONALLY, THAT'S FOR SURE

WHAT DO YOU THINK I AM, STUPID?

\vspace*{\fill}

\newpage


\titledquestion{Checking Every Element}

\textbf{Checking Every Element}

For the type of polymorphic trees, consider the functions \code{treeAll}, which checks
that every element of a tree satisfies a predicate, and \code{all}, which checks
that every element of a list satisfies a predicate. Additionally provided is the
\code{inorder} function, which returns the inorder traversal of a tree.
\begin{codeblock}
  datatype 'a tree = Empty | Node of 'a tree * 'a * 'a tree

  fun treeAll p Empty = true
    | treeAll p (Node (l, x, r)) =
        p x andalso treeAll p l andalso treeAll p r

  fun all p [] = true
    | all p (x :: xs) = p x andalso all p xs

  fun inorder Empty = []
    | inorder (Node (l, x, r)) = inorder l @ (x :: inorder r)
\end{codeblock}

Checking a predicate over a tree directly should agree with flattening to the
inorder list and checking that. Prove it.

\begin{parts}
  \part[15]\hfill

    \textbf{Theorem.} For all values \code{p : 'a -> bool} and \code{T : 'a tree}:
    $$\code{treeAll p T} \;\eeq\; \code{all p (inorder T)}$$

    You may use the following facts without proof:

    \textbf{Lemma 1:} For any value
    \code{T : 'a tree}, \code{inorder T} evaluates to a value.

    \textbf{Lemma 2:} For total \code{p} and any
    value \code{L : 'a list}, \code{all p L} evaluates to a value.

    \textbf{Lemma 3:} For total \code{p} and any
    value \code{T : 'a tree}, \code{treeAll p T} evaluates to a value.

    \textbf{Lemma 4:} For total \code{p} and any
    values \code{A, B : 'a list},
    $$\code{all p (A @ B)} \;\eeq\;
      \code{all p A andalso all p B}$$

    \begin{constraint}
      You may assume \code{p} is total.
    \end{constraint}

    \begin{solutionorbox}[50em]

      By structural induction on \code{T : 'a tree}.

      \vspace{6pt}
      \textbf{Base case:} \code{T = Empty}.
      \begin{align*}
        \code{treeAll p Empty}
        &\eeq \code{true} \tag{clause 1 of \code{treeAll}} \\
        &\eeq \code{all p []} \tag{clause 1 of \code{all}} \\
        &\eeq \code{all p (inorder Empty)} \tag{clause 1 of \code{inorder}}
      \end{align*}

      \vspace{6pt}
      \textbf{Inductive case:} \code{T = Node (l, x, r).}

      \noindent IH$_l$: \code{treeAll p l} $\eeq$ \code{all p (inorder l)}.

      \noindent IH$_r$: \code{treeAll p r} $\eeq$ \code{all p (inorder r)}.

      \noindent WTS: \code{treeAll p (Node (l,x,r))} $\eeq$
      \code{all p (inorder (Node (l,x,r)))}.

      \vspace{4pt}
      We start from the \textbf{right}-hand side and reach the left:
      \begin{align*}
        & \code{all p (inorder (Node (l,x,r)))} \\
        &\eeq \code{all p (inorder l @ (x :: inorder r))}
          \tag{clause 2 of \code{inorder}} \\
        &\eeq \code{all p (inorder l)} \\
        & \quad\; \code{andalso all p (x :: inorder r)}
          \tag{Lemmas 1, 4} \\
        &\eeq \code{all p (inorder l)} \\
        & \quad\; \code{andalso (p x andalso all p (inorder r))}
          \tag{clause 2 of \code{all}} \\
        &\eeq \code{p x andalso all p (inorder l)} \\
        & \quad\; \code{andalso all p (inorder r)}
          \tag{\code{andalso} assoc.\ / comm., Lemma 2 and \code{p} total} \\
        &\eeq \code{p x andalso treeAll p l andalso treeAll p r}
          \tag{IH$_l$, IH$_r$, Lemma 3} \\
        &\eeq \code{treeAll p (Node (l,x,r))}
          \tag{clause 2 of \code{treeAll}}
      \end{align*}

      which is the left-hand side. By the principle of structural induction, the
      theorem holds for all \code{T : 'a tree}. $\qed$
    \end{solutionorbox}

    \newpage

    \begin{solutionorbox}[50em]

    \end{solutionorbox}
\end{parts}
